% UTF8 表示使用通用国际字符作为内码，ctexart 表示这是一篇中文论文
\documentclass[UTF8]{ctexart} 
% 通过调用宏包 geometry 调整页面设置为 A4 纸，页边距 2.6 cm
\usepackage[a4paper,margin=2.6cm]{geometry}

\usepackage{amsmath, amssymb, amsthm} % 数学公式、符号、定理
\usepackage{graphicx}                 % 插图
\usepackage[colorlinks=true,linkcolor=blue]{hyperref} % 超链接与交叉引用

% 设置文章标题
\title{复变函数学习笔记：3.2节“留数公式”}
% 设置撰写日期
\date{}

% 设置环境 theorem (这是用户自定义环境名，相当于一个变量)
\newtheorem{theorem}{定理}[section]
% 设置引理环境，直接套用theorem，后果是和定理共享编号
\newtheorem{lemma}[theorem]{引理}
% 设置定义环境，和定理共享编号
\newtheorem{definition}[theorem]{定义}
% 设置注记环境，独立编号，和节绑定
\newtheorem{remark}{注记}[section]
% 让公式编号的计数器前面自动加上节，比如式 (2.1)
\numberwithin{equation}{section}

% 自定义命令，专门用于实数域和复数域
\newcommand{\RR}{\mathbb{R}}
\newcommand{\CC}{\mathbb{C}}

\begin{document}

\maketitle

\section{逻辑脉络 (Logical Flow)}
本节是复变函数中最核心、最实用的部分之一。其逻辑脉络承前启后，将 3.1 节中纯局部的概念（留数）与第 2 章中的全局积分工具（柯西定理）完美结合，并最终落地于解决实际的微积分问题。具体脉络如下：
\begin{enumerate}
    \item \textbf{单极点圆周积分}：首先解决最基本的情形——如果积分路径是一个圆，且内部只有一个极点，积分结果是什么？利用“钥匙孔”（keyhole）围道，将问题转化为绕极点的小圆积分，从而提取出留数。
    \item \textbf{多极点与一般围道推广}：从单极点圆周推广到多极点圆周，再进一步推广到包含多个极点的“玩具围道”（Toy contour，如矩形、半圆等）。这确立了著名的\textbf{留数公式}：闭曲线积分完全由其内部包围的所有极点的留数之和决定。
    \item \textbf{方法落地：计算实积分}：留数公式的最强大应用在于计算广义黎曼积分（如 $\int_{-\infty}^\infty f(x)dx$）。教材在 3.2.1 小节中通过三个递进的经典例子，展示了如何通过“解析延拓 $\to$ 构造闭合围道 $\to$ 控制边界积分 $\to$ 计算内部留数”的标准化流程，将困难的实积分转化为简单的代数计算。
\end{enumerate}

\section{主要内容梳理 (Main Content)}

\subsection{留数公式 (The Residue Formula)}
教材首先给出了单一极点在圆周内的情况，并随后将其推广到最一般的形式：

\begin{theorem}[留数公式]
假设 $f$ 在包含“玩具围道”（Toy contour）$\gamma$ 及其内部的开集上全纯，除了在 $\gamma$ 内部的有限个点 $z_1, z_2, \dots, z_N$ 处有极点。假定 $\gamma$ 具有正定向（逆时针），那么：
$$ \int_\gamma f(z) dz = 2\pi i \sum_{k=1}^N \mathrm{res}_{z_k} f $$
\end{theorem}

\begin{remark}
\textbf{证明的核心思想}：通过构造“多重钥匙孔”围道（multiple keyhole），利用柯西定理将沿着大围道 $\gamma$ 的积分，转化为沿着环绕各个极点 $z_k$ 的小圆 $C_k$ 的积分之和。对于每个小圆 $C_k$，函数 $f(z)$ 可写为全纯部分 $G(z)$ 与主部分（Principal part）之和，而全纯部分与所有 $(z-z_k)^{-m} (m > 1)$ 的项积分均为 0，唯独剩下 $(z-z_k)^{-1}$ 的系数（即留数），每次积出 $2\pi i \cdot \mathrm{res}_{z_k} f$。
\end{remark}

\subsection{留数公式的应用：计算实积分}
留数微积分（Calculus of residues）提供了一种计算各种定积分的强大技术。核心思想是将定义在实轴上的函数 $f(x)$ 扩展为复平面上的函数 $f(z)$，然后选择一个包含实轴某一段 $[-R, R]$ 的合适的闭合玩具围道 $\gamma_R$，使得当 $R \to \infty$ 时，其余边界上的积分趋于零或容易处理。

教材通过三个经典例子展示了不同的围道选择策略：

\paragraph{例1：半圆围道 (Semicircle Contour)}
计算积分：
$$ \int_{-\infty}^\infty \frac{dx}{1+x^2} = \pi $$
\textbf{策略}：考虑复函数 $f(z) = 1/(1+z^2)$，它在 $z = \pm i$ 处有简单极点。构造由实轴线段 $[-R, R]$ 和上半平面的大半圆 $C_R^+$ 组成的闭合围道 $\gamma_R$。由于上半平面内只有极点 $z=i$，其留数为 $1/(2i)$。当 $R \to \infty$ 时，由于分母是 $z^2$ 级别，大半圆上的积分被压制为 $0$。因此实积分直接等于 $2\pi i \cdot (1/2i) = \pi$。

\paragraph{例2：矩形围道 (Rectangle Contour)}
计算积分（其中 $0 < a < 1$）：
$$ \int_{-\infty}^\infty \frac{e^{ax}}{1+e^x} dx = \frac{\pi}{\sin \pi a} $$
\textbf{策略}：对于含有指数函数的被积函数，通常选择\textbf{矩形围道}。考虑 $f(z) = e^{az}/(1+e^z)$，极点位于分母为零处，即 $z = \pi i$。取顶点为 $-R, R, R+2\pi i, -R+2\pi i$ 的矩形。因为 $e^z$ 的周期性为 $2\pi i$，矩形上下两边的积分可以合并（差一个常数因子 $e^{2\pi i a}$），而左右两竖边的积分在 $R \to \infty$ 时趋于零，从而通过留数求出原积分。

\paragraph{例3：带振荡因子的矩形围道}
证明双曲正割函数的一个变体是其自身的傅里叶变换：
$$ \int_{-\infty}^\infty \frac{e^{-2\pi i x \xi}}{\cosh \pi x} dx = \frac{1}{\cosh \pi \xi} $$
\textbf{策略}：与例2类似，利用被积函数中 $\cosh(\pi z)$ 的周期特性，同样选取适当高度的矩形围道。此例中内部包含两个简单极点 $\alpha = i/2$ 和 $\beta = 3i/2$。计算它们的留数并处理矩形顶部路径带来的振荡与衰减，最后合并得到结果。这为后续章节（如 Theta 函数、双平方和定理）中研究此函数奠定了基础。

\end{document}